%% kapa: exam=1 round=32 year=2021 period=2 form=A subject=appraisal_law subject_ko=감정평가관계법규 seq=8 number=8
%% source: Q-Net 감정평가사 기출문제 (cst00309) artlSeq=5208397 file=extracted/1008/5208397/A형/2021년도 제32회 감정평가사 1차 2교시 A형.pdf
%% conversion: kordoc 4.24.51 -> tools/gen1.py (auto); review=proofread-sonnet; flags=none
\begin{question}{8}
국토의 계획 및 이용에 관한 법령상 지구단위계획에 관한 설명이다. (\quad)에 들어갈 내용으로 각각 옳은 것은?
\begin{framedtext}
주민의 입안제안에 따른 지구단위계획에 관한 (ㄱ)결정의 고시일부터 (ㄴ) 이내에 이 법 또는 다른 법률에 따라 허가·인가·승인 등을 받아 사업이나 공사에 착수하지 아니하면 그 (ㄴ)이 된 날의 다음날에 그 지구단위계획에 관한 (ㄱ)결정은 효력을 잃는다.
\end{framedtext}
\begin{choices}
  \item ㄱ: 도시·군기본계획, ㄴ: 3년
  \item ㄱ: 도시·군기본계획, ㄴ: 5년
  \item ㄱ: 도시·군관리계획, ㄴ: 1년
  \item ㄱ: 도시·군관리계획, ㄴ: 3년
  \item ㄱ: 도시·군관리계획, ㄴ: 5년
\end{choices}
\end{question}
\answer{5}
